% KI-Verzeichnis nach dem Leitfaden im Nachspann hinter dem Literaturverzeichnis.
% Kein Prompt-Verzeichnis (Vorgabe des Betreuers im Milestone-Gespräch vom 4. November 2026, hat Vorrang vor dem Leitfaden).
% Verfasser: Modellversionen und Zeitraum ggf. ergänzen, z. B. "Claude Opus 4.x, Juli bis November 2026".
\section*{KI-Verzeichnis}
\addcontentsline{toc}{section}{KI-Verzeichnis}

Anthropic. (2026). Claude Opus [Großes Sprachmodell]. https://claude.ai

Anthropic. (2026). Claude Code [KI-Programmierassistent]. https://claude.com/claude-code

\medskip
\noindent Fragestellung, Analyse, Entwurfsentscheidungen, Messungen und Schlussfolgerungen sind eigene Arbeit. Die
Werkzeuge halfen beim Programmieren, überarbeiteten den selbst geschriebenen Text und unterstützten die Prüfung von Quellen und
Zwischenständen. Sie waren in keinem Fall Quelle einer Aussage. Jede Aussage der Arbeit stützt sich auf eine im
Literaturverzeichnis genannte Quelle, eine eigene Messung oder den Quellcode des Prototyps. Personenbezogene Daten
wurden nicht eingegeben. Die Vorgaben der FOM, die Leitlinien der Deutschen Forschungsgemeinschaft zur Sicherung guter
wissenschaftlicher Praxis (Kodex 2023) und die Datenschutz-Grundverordnung wurden eingehalten.

\begin{table}[H]
  \caption{Eingesetzte KI-Werkzeuge nach Aufgabe}
  \label{tab:ki-verzeichnis}
  \small
  \begin{tabularx}{\textwidth}{@{}>{\RaggedRight\arraybackslash}p{0.20\textwidth} >{\RaggedRight\arraybackslash}X >{\RaggedRight\arraybackslash}p{0.36\textwidth}@{}}
    \hline
    Aufgabe & Umfang & Eigener Anteil und Prüfung \\
    \hline
    Programmierung & Unterstützung beim Schreiben des Codes, vor allem beim Bootstrapping, also beim Grundgerüst der
      Projekte & Architektur, Schnittstellen und Sicherheitsvorgaben selbst festgelegt, erzeugten Code geprüft,
      angepasst und mit Prüffällen, Differenztest und Lastläufen getestet \\
    Überarbeitung des Textes & Kürzen, Umformulieren und Umstellen selbst geschriebener Abschnitte, nach
      Gutachten auch Neufassung einzelner Absätze aus eigenen Inhalten, Rechtschreibung, Konsistenz von Zahlen und
      Verweisen & Inhalt, Gliederung, Titel und Forschungsfragen selbst erarbeitet, jede Änderung einzeln angenommen,
      angepasst oder verworfen \\
    Literaturprüfung & Suche nach Kandidaten für Quellen, Abgleich von Zitaten und Seitenangaben mit den Volltexten,
      Ergänzung bibliografischer Angaben & Auswahl und Lektüre jeder zitierten Quelle, nicht auffindbare Vorschläge
      verworfen \\
    Abbildungen & Umsetzung in \LaTeX{} und Diagrammcode nach eigenen Vorgaben & Inhalt aus eigenem Entwurf und
      eigenen Messdaten, Ergebnis gegen beide geprüft \\
    Rückmeldung zu Zwischenständen & Simulierte Gutachten aus mehreren Perspektiven & Jeden Befund bewertet und über
      Übernahme oder Ablehnung selbst entschieden \\
    \hline
  \end{tabularx}
\end{table}
